\section{Common level labels and their Heisenberg phases}
\label{sec:sr-level-phase-proof}

We verify the finite test-vector norm and overlap in
\eqref{eq:sr-test-gram} using the actual shifted level lattices.
The same calculation supplies one multiplier for all auxiliary
positive definite lattices.

Let \(L_0\) be the fixed polarized \(\mathcal O_E\)-lattice of
rank \(r\), selfdual for its rational alternating trace form
\(B\).  Thus \(B(av,z)=B(v,\overline a z)\).
Let \(\mathfrak a\) be the finite level ideal, with
\[
 (\mathfrak a,\overline{\mathfrak a})=1,\qquad
 (\mathfrak a,2)=1,\qquad M\in\mathfrak a,\quad M\ne0.
\]
Choose \(w\in\mathcal O_E\) with \(w+\overline w=1\).
All Heisenberg operators and all vacua below use this one \(w\).
Put \(\mathfrak j=\mathfrak a\overline{\mathfrak a}^{-1}\).
The lattices \(L\) under consideration are selfdual and agree
with \(L_0\) after localization at every prime dividing
\(\mathfrak a\overline{\mathfrak a}\).
This is equivalent to having a common denominator, prime to
those ideals, carrying each lattice into the other.

The auxiliary lattices all lie in one fixed Hermitian space,
and this is essential for the projections below.  More explicitly,
fix
\[
 V_+=E^n,\qquad
 h_+(v,v')=\sum_{i=1}^n v_i\overline{v'_i},\qquad
 \Lambda_0=\mathcal O_E^n.
\]
Assume \(n\) is even.  Every \(\Lambda\) below is an
\(\mathcal O_E\)-submodule of this \(V_+\), selfdual for this
one \(h_+\), and standard at the primes of \(M\).
The form is positive definite at each complex embedding.
The lattice \(\Lambda_0\) is selfdual: pairing with the coordinate
vectors shows that its Hermitian dual has all coordinates in
\(\mathcal O_E\), and the reverse inclusion follows by summing
products of integral coordinates.  No individual Hermitian form
is assigned to \(\Lambda\).

Put \(V=E^r\), with its fixed polarized trace form \(B\).
There is one rational symplectic space
\[
 \mathbb W=V_+\otimes_E V
        \simeq \operatorname{Mat}_{n,r}(E),\qquad
 B_n(Y,Z)=\sum_{i=1}^n B(Y_i,Z_i).
\]
On pure tensors its pairing is
\(B_n(v\otimes u,v'\otimes u')
  =B(h_+(v,v')u,u')\), by the adjointness of multiplication
by \(E\).  Thus the perfect tensor pairings used below are
restrictions of this same \(B_n\).
Choose a rational polarization of \(\mathbb W\) once, for example
by choosing a symplectic integral basis of \(L_0\), and take
\[
 {\cal H}=L^2(\mathbb A_f^{ng}),\qquad
 {\cal S}={\cal S}(\mathbb A_f^{ng}),\qquad
                   2g=\dim_{\mathbb Q}V.
\]
Use the standard rational additive character and its selfdual
Haar measure.  In the rational-function notation of the source,
\({\cal S}\) is the space of rational Schwartz--Bruhat functions
on \(\operatorname{Mat}_{n,g}(\mathbb Q)\): restriction gives the
identification, because \(\mathbb Q^{ng}\) is dense in
\(\mathbb A_f^{ng}\).  A common support denominator and an open
period subgroup make the inverse extension well defined; its
finite residue average is the same \(L^2\) inner product.
All vacua, projections and translated vectors in this section
are vectors in this one \({\cal H}\).

The group \(U(V,B)\) acts on \(\mathbb W\) on each of the \(n\)
rows, preserving \(B_n\).  We use one genuine unitary Weil
action \(W\) of this group on \({\cal S}\), with the covariance
convention
\[
 W_\gamma S(Y)W_\gamma^{-1}=S(Y\gamma^t).
\]
Here genuine means \(W_{\gamma'\gamma}=W_{\gamma'}W_\gamma\).
This action is fixed before the lattices are varied.
The standard metaplectic Schr\"odinger representation supplies
it: the action on \(n\) rows is the \(n\)-fold tensor power of
the action for \(V\); the two lifts of a symplectic map differ
by \(-1\), acting by \((-1)^n=1\) on that tensor power.
It therefore descends to the symplectic group when \(n\) is
even.  Fourier transform with selfdual measure, quadratic
characters of modulus one and normalized linear changes of
variables give its unitary operators and preserve \({\cal S}\).
This is the standard metaplectic background, with the additive
character, measure and parity now specified; see Weil,
\emph{Sur certains groupes d'op\'erateurs unitaires},
Acta Math.\ 111 (1964), 143--211,
\url{https://doi.org/10.1007/BF02391012}.

The theta-normalized action in the source is characterized by
\[
 \operatorname{thA}(W_\gamma\phi;\gamma x,z)
                         =\operatorname{thA}(\phi;x,z).
\]
The usual Poisson theta transformation fixes this scalar
normalization.  Its uniqueness, and hence its genuine group law,
can also be checked directly.  At one frame \(x\), the invertible
period-coordinate change and the nonvanishing Gaussian factor
reduce equality of these entire functions of \(z\) to equality of
\[
 \sum_{a\in\mathbb Q^{ng}}
       \phi(a)\exp\{\pi i\,\operatorname{tr}(a\tau a^t)
                              +2\pi i\,\operatorname{tr}(az^t)\},
 \qquad \operatorname{Im}\tau>0 .
\]
Put the two test functions over one support denominator.
Absolute Gaussian convergence permits Fourier coefficient
extraction on a real period torus; the coefficient at \(a\) is
its displayed nonzero Gaussian times \(\phi(a)\).
Thus equality gives equality of the test functions.
Applying the theta transformation twice proves the group law
by this uniqueness, without choosing phases separately for
different lattices.  In the frozen archive the precise producer
is \srcid{WeilActionPreservesThetaAndInnerProduct.wG\_laws};
its inputs are the fixed symplectic basis and evenness of \(n\),
and its conclusions include unitarity, linearity, the group law
and the above Heisenberg covariance.

With this fixed ambient model, set
\[
 A_{\Lambda,L}=\Lambda\otimes_{\mathcal O_E}\mathfrak jL,\quad
 H_{\Lambda,L}=\Lambda\otimes_{\mathcal O_E}
                         2\overline{\mathfrak a}L,\quad
 K_{\Lambda,L}=\Lambda\otimes_{\mathcal O_E}
                         2\overline{\mathfrak a}\mathfrak aL.
\]
The first lattice is selfdual.
Indeed tensoring the perfect Hermitian pairing on \(\Lambda\)
with the perfect alternating pairing on \(L\) is perfect;
and the twist by \(\mathfrak j\) preserves selfduality because
\(\mathfrak j\overline{\mathfrak j}=1\).
These statements can be checked over every local DVR and then
intersected inside the rational space.

Fix the Schr\"odinger operators with convention
\[
 S(u)S(v)=e^{-\pi iB_n(u,v)}S(u+v),\qquad
 S(u)S(v)=e^{-2\pi iB_n(u,v)}S(v)S(u).
\]
For a rational vector \(\delta\), set
\begin{equation}\label{eq:sr-common-eigenoperator}
 \epsilon_w(\delta)=e^{\pi iB_n(\delta,w\delta)},\qquad
 U_\delta=\epsilon_w(\delta)S(\delta).
\end{equation}
On any integral \(\mathcal O_E\)-stable lattice \(A\), these
operators form an additive representation.
In fact the quotient of \(U_uU_v\) by \(U_{u+v}\) is
\[
 \exp\{-\pi i[
 B_n(u,wv)+B_n(v,wu)+B_n(u,v)]\}
       =\exp\{-2\pi iB_n(u,wv)\}=1.
\]
Here adjointness and \(w+\overline w=1\) give the equality,
and integrality gives the final value.

There is a one-dimensional vacuum line fixed by all
\(U_\delta\), \(\delta\in A\), for every selfdual \(A\).
For completeness, choose a symplectic integral basis of one
selfdual lattice and its Schr\"odinger realization.
The standard compact-open indicator, translated and multiplied
by half-integral characters to account for the signs
\(\epsilon_w\) on this basis, gives its vacuum.
For a second commensurable selfdual lattice, reduce to the
finite Heisenberg quotient of their intersection.
In the model of functions on a finite maximal isotropic group,
translations and the distinct dual characters have each joint
eigenline of dimension one.
They give the corresponding vacuum and show uniqueness.
The construction uses finite sums of rational Heisenberg
translations, so the vacuum remains Schwartz--Bruhat.

The phases can be chosen compatibly with the prescribed genuine
unitary Weil action.  First choose unit vacua \(v_{0,L}\) for
\(A_{0,L}=\mathcal O_E^n\otimes\mathfrak jL\), once for each \(L\).
For a level-group element \(\gamma\), write
\[
 W_\gamma v_{0,L}=\mu(\gamma,L)v_{0,\gamma L}.
\]
The lines have just been proved one-dimensional, so the scalar
exists, has modulus one, and satisfies the cocycle relation
forced by the same genuine action \(W\).
For an arbitrary auxiliary \(\Lambda\), let \(P_{\Lambda,L}\)
be the orthogonal projector onto the
\(A_{\Lambda,L}\)-vacuum line, obtained by averaging
\(U_\delta\) over its finite orbit on \(v_{0,L}\), and set
\[
 v_{\Lambda,L}=
  \frac{P_{\Lambda,L}v_{0,L}}
       {\|P_{\Lambda,L}v_{0,L}\|}.
\]
The projection is nonzero: the finite Heisenberg calculation
of Lemma~\ref{lem:sr-selfdual-count}'s subsequent vacuum argument
gives overlap of modulus the inverse square root of the
relative index, since the semicharacters agree on the
intersection.
Weil covariance intertwines the two projectors.
Unitarity preserves their positive normalization denominators,
and hence
\begin{equation}\label{eq:sr-common-vacuum-multiplier}
 W_\gamma v_{\Lambda,L}
       =\mu(\gamma,L)v_{\Lambda,\gamma L}.
\end{equation}
In particular the multiplier is the same for every \(\Lambda\);
it is not chosen independently for each auxiliary theta vector.

Choose \(r\) vectors \(\beta_1,\ldots,\beta_r\in L_0\)
spanning \(L_0/\mathfrak aL_0\).
They form a basis over \(\mathcal O_E/\mathfrak a\):
local freeness gives cardinality
\(\operatorname N(\mathfrak a)^r\), so the surjection from the
free module of that same cardinality is bijective.
Let
\(J_{\rm lev}\subset\operatorname{GL}_r(\mathcal O_E/\mathfrak a)\)
be the prescribed finite subgroup, with character
\(\lambda:J_{\rm lev}\to\mathbb C^\times\).
Its values have modulus one, since \(J_{\rm lev}\) is finite.
For \(k\in J_{\rm lev}\), prescribe the residue of a row tuple by
\[
 \bar y_i(k)=\sum_j k_{ji}\bar\beta_j\quad(1\le i\le r),
 \qquad \bar y_i(k)=0\quad(r<i\le n).
\]
These are distinct labels, indexed exactly by \(J_{\rm lev}\).

There is a canonical residue isomorphism
\begin{equation}\label{eq:sr-level-quotient}
 H_{\Lambda,L}/K_{\Lambda,L}
              \simeq (L_0/\mathfrak aL_0)^n.
\end{equation}
To prove it, localize at a prime \(\mathfrak p\mid\mathfrak a\).
The elements \(2\) and
\(\overline{\mathfrak a}\) are units there,
\(\Lambda\) is standard there because \(M\in\mathfrak a\),
and \(L_{\mathfrak p}=(L_0)_{\mathfrak p}\).
Thus the localized quotient is precisely the usual reduction
of \(L_0^n\) modulo \(\mathfrak a\).
At primes outside \(\mathfrak a\) the quotient is zero.
The local identifications combine by the Chinese remainder
theorem, proving \eqref{eq:sr-level-quotient}.

Choose any representatives \(y_k\in H_{\Lambda,L}\) of the
prescribed labels and define the actual test vector
\begin{equation}\label{eq:sr-actual-level-test}
 \Phi_{\Lambda,L}
    =\sum_{k\in J_{\rm lev}}\lambda(k)S(y_k)v_{\Lambda,L}.
\end{equation}
This definition is independent of the representatives, with no
missing root-of-unity factor.
Indeed for \(h\in K_{\Lambda,L}\), one has
\(h\in A_{\Lambda,L}\), because the ideal ratio is
\(2\overline{\mathfrak a}^{\,2}\subset\mathcal O_E\).
Also \(y_k,h\in2(\Lambda\otimes L)\).
This untwisted tensor lattice is integral and stable under \(w\),
so
\[
 B_n(y_k,h)\in4\mathbb Z,\qquad
 B_n(h,wh)\in4\mathbb Z .
\]
Thus \(\epsilon_w(h)=1\), \(S(h)v_{\Lambda,L}=v_{\Lambda,L}\),
and the Heisenberg product law gives
\(S(y_k+h)v_{\Lambda,L}=S(y_k)v_{\Lambda,L}\).
The explicit factor \(2\) in the two quotient lattices is what
makes this representative change exact, including at two.

We can now use common representatives for two lattices \(L,L'\).
Their modules \(H_{\Lambda,L}\) and \(H_{\Lambda,L'}\) have the
same localization at primes of \(\mathfrak a\).
Their intersection therefore still surjects onto
\((L_0/\mathfrak aL_0)^n\), by the same local argument and CRT.
Choose \(y_k\) in this intersection.
The just-proved independence means that these simultaneous
representatives compute both test vectors in
\eqref{eq:sr-actual-level-test}.

\begin{srproposition}[The common operator separates distinct labels]
\label{prop:sr-level-label-orthogonality}
With those representatives,
\[
 \langle S(y_k)v_{\Lambda,L},
              S(y_{k'})v_{\Lambda,L'}\rangle=0
                      \qquad(k\ne k').
\]
\end{srproposition}

\begin{proof}
Put \(A=A_{\Lambda,L}\), \(A'=A_{\Lambda,L'}\).
For every \(\delta\in A\cap A'\), the very same operator
\(U_\delta\) from \eqref{eq:sr-common-eigenoperator} acts on the
two vectors with eigenvalues
\[
 e^{-2\pi iB_n(\delta,y_k)},\qquad
 e^{-2\pi iB_n(\delta,y_{k'})},
\]
respectively.  This follows directly by commuting \(S(\delta)\)
past \(S(y)\); the common factor \(\epsilon_w(\delta)\)
cancels against its eigenvalue on each vacuum.
If their inner product were nonzero, unitarity would force
these eigenvalues to agree.  Thus
\[
 B_n(\delta,y_k-y_{k'})\in\mathbb Z
                           \quad(\delta\in A\cap A').
\]
Consequently
\[
 y_k-y_{k'}\in(A\cap A')^\#=A^\#+(A')^\#=A+A'.
\]
The dual-intersection identity can be verified without a
duality convention: the exact sequence
\(0\to A\cap A'\to A\oplus A'\to A+A'\to0\)
splits as abelian groups since \(A+A'\) is free.
An integral functional on \(A\cap A'\) therefore extends to
one on \(A\oplus A'\), giving the asserted sum of the duals.

At \(\mathfrak p\mid\mathfrak a\), both \(A\) and \(A'\)
localize to
\(\mathfrak a_{\mathfrak p}(L_0)_{\mathfrak p}^{\,n}\):
the conjugate ideal is a unit there, and the auxiliary lattice
is standard there.
Thus \(y_k-y_{k'}\) has zero residue in
\((L_0/\mathfrak aL_0)^n\).
By \eqref{eq:sr-level-quotient} the two labels coincide.
The \(\beta_j\)'s are a residue basis, so \(k=k'\), a
contradiction.
\end{proof}

\begin{srproposition}[Exact level Gram normalization]
\label{prop:sr-exact-level-gram}
For the single construction above, with
\(b=|J_{\rm lev}|\),
\[
 \langle\Phi_{\Lambda,L},\Phi_{\Lambda,L'}\rangle
         =b\,\langle v_{\Lambda,L},v_{\Lambda,L'}\rangle.
\]
In particular
\[
 \|\Phi_{\Lambda,L}\|^2=b,\qquad
 |\langle\Phi_{\Lambda,L},\Phi_{\Lambda,L'}\rangle|
      \le b\,m(L,L')^{-n/2}.
\]
These are \eqref{eq:sr-test-gram}, now with the level phases fixed.
\end{srproposition}

\begin{proof}
Expand the two finite sums using their common representatives.
Proposition~\ref{prop:sr-level-label-orthogonality} removes all
terms with different labels.
For equal labels, the common \(S(y_k)\) is unitary and
\(|\lambda(k)|=1\), so the summand is exactly
\(\langle v_{\Lambda,L},v_{\Lambda,L'}\rangle\).

To compute the index in the vacuum overlap, note that
\(\Lambda\) is projective of rank \(n\), and the twist
\(\mathfrak j\) is an invertible module.
Local bases therefore give
\[
 A_{\Lambda,L}\cap A_{\Lambda,L'}
       =\Lambda\otimes\mathfrak j(L\cap L'),\qquad
 [A_{\Lambda,L}:A_{\Lambda,L}\cap A_{\Lambda,L'}]
       =[L:L\cap L']^n.
\]
For the index assertion, decompose each finite local quotient
into cyclic modules over its local DVR; tensoring a rank-\(n\)
free module makes \(n\) copies, and tensoring an invertible
ideal preserves its length.  Multiplication of the local
cardinalities proves the formula.
The normalized finite Heisenberg overlap is the inverse
square root of this index, as proved in
Section~\ref{sec:sr-theta-gram}.
This gives the stated bound and the diagonal norm.
\end{proof}

The multiplier on the actual level test vectors includes the fixed
level character.  If the residue action of \(\gamma\) on the
\(\beta\)-basis is \(k_\gamma\in J_{\rm lev}\), applying
\(\gamma\) sends the label \(k\) to \(k_\gamma k\).
Reindexing the finite sum, with the representative-independence
already proved, gives the exact formula
\[
 W_\gamma\Phi_{\Lambda,L}
   =\widetilde\mu(\gamma,L)\Phi_{\Lambda,\gamma L},
 \qquad
 \widetilde\mu(\gamma,L)
       =\mu(\gamma,L)\lambda(k_\gamma)^{-1}.
\]
Thus it is \(\widetilde\mu\) that enters the lattice Gram
automorphy identity and the inverse law for holomorphic families.
It remains independent of \(\Lambda\), and has modulus one.
This fixes the possible level-character normalization shift.

All choices in this section are made in the fixed standard
Hermitian space \(V_+\), the fixed rational symplectic space
\(\mathbb W\) and its single Schr\"odinger Hilbert space, from
one trace form, one semicharacter \(w\), one genuine Weil action,
the fixed level character, and common residue representatives.
The projection construction proves
\eqref{eq:sr-common-vacuum-multiplier} with a multiplier
independent of the auxiliary lattice.
It does not construct the arithmetic polarized lattice or its
integral model; those remain part of the family construction.
In the frozen source the displayed phase calculation is
\srcid{ThetaTestFunctionNormAndOverlapBound.pair_int};
the local residue implication is \srcid{cong_of_pair};
and the final diagonal sum is \srcid{phi_norm_overlap}.


The fixed-model assertion here agrees literally with the source
types.  In \srcid{FirstCoordinateAdmissible}, \srcid{Lt E n}
is a submodule of \(E^n\), and \srcid{Vset E n M} uses
\srcid{herm 1} for every auxiliary lattice.
In \srcid{ThShort}, \srcid{GenLat} has only the submodule
field \srcid{La}, and \srcid{GenLat.Unimod} uses the same
coordinate Hermitian pairing.  In \srcid{IsSystem},
\srcid{TF n g} is fixed independently of \(\Lambda\).
These are lattice variations inside a fixed space, so no
unproved isometric identification between varying rational
Hermitian forms is being used in \(P_{\Lambda,L}v_{0,L}\).
The arithmetic polarized model and the compatibility of the
global doubled theta kernel with its integral geometric
realization remain separate requirements.

\subsection{Zero overlap outside the prescribed level group}
\label{subsec:sr-level-support}

Continue in the single fixed Hermitian and Schr\"odinger model
of Section~\ref{sec:sr-level-phase-proof}.
Define the level group \(\Delta\) to consist of those rational
unitary maps \(\gamma\) which preserve \(L_0\) locally at the
primes of \(\mathfrak a\) and act on
\(L_0/\mathfrak aL_0\), in the \(\beta\)-basis, by a matrix
\(k_\gamma\in J_{\rm lev}\).
This is a subgroup: preservation and the residue matrices
are stable under multiplication and inverse.
Equivalently it is the subgroup generated by the rational
unitary maps satisfying this condition.
The latter is the source's definition \srcid{Level.grp}:
its \srcid{Cong} relation is equality in these local residue
modules, since a denominator prime to \(\mathfrak a\) is
invertible there.  Conversely, local residue equality can be
multiplied by a common denominator prime to
\(\mathfrak a\), giving exactly that \srcid{Cong} relation.

\begin{srproposition}\label{prop:sr-level-outside-zero}
For every rational unitary \(\gamma\notin\Delta\), every two
allowed \(L,L'\), and every allowed \(\Lambda\),
\[
       \langle\Phi_{\Lambda,L},
                        W_\gamma\Phi_{\Lambda,L'}\rangle=0 .
\]
\end{srproposition}

\begin{proof}
Suppose the inner product is nonzero.
Expand the two finite labelled sums and use Weil covariance.
Some pair of labels \(k,k'\) therefore satisfies
\[
 \langle S(y_k)v_{\Lambda,L},
              S(\gamma y'_{k'})W_\gamma v_{\Lambda,L'}\rangle
                                                        \ne0 .
\]
We regard rows of a tuple as column vectors when applying
\(\gamma\).
Put \(A=A_{\Lambda,L}\) and \(A'=\gamma A_{\Lambda,L'}\).
Both are selfdual lattices for the same \(B_n\).
The vector \(W_\gamma v_{\Lambda,L'}\) is fixed by
\(U_\delta\), \(\delta\in A'\): covariance gives this, and
\[
 B_n(\gamma z,w\gamma z)=B_n(z,wz)
\]
since \(\gamma\) is \(E\)-linear and unitary.
Consequently the common operators \(U_\delta\),
\(\delta\in A\cap A'\), have equal eigenvalues on the two
vectors with nonzero inner product.
Exactly as in
Proposition~\ref{prop:sr-level-label-orthogonality}, this yields
\begin{equation}\label{eq:sr-level-overlap-difference}
              y_k-\gamma y'_{k'}\in A+A'.
\end{equation}

Fix a prime \(\mathfrak q\mid\mathfrak a\), and put
\(e=v_{\mathfrak q}(\mathfrak a)>0\).
At this prime, \(\Lambda\) is standard,
\(L=L'=L_0\), and
\(\mathfrak j=\mathfrak a\overline{\mathfrak a}^{-1}\)
is \(\mathfrak q^e\) times a unit.
Choose a local basis of \(L_0\), and let \(U,U'\) be the
\(r\)-by-\(r\) matrices whose columns are the first \(r\)
rows of \(y_k,y'_{k'}\), in this basis.
These columns are integral and give bases modulo
\(\mathfrak q^e\), because their labels come from invertible
\(k,k'\) and the residue basis \(\beta\).
Their determinants are therefore units, so
\(U,U'\in\operatorname{GL}_r(\mathcal O_{E,\mathfrak q})\).

Equation~\eqref{eq:sr-level-overlap-difference} gives integral
matrices \(R,S\) with
\[
                 U-\gamma U'=\pi^e R+\pi^e\gamma S ,
\]
where \(\pi\) is any local uniformizer; the unit in the
fractional ideal has been absorbed in \(R,S\).
Thus
\[
       \gamma=(U-\pi^eR)(U'+\pi^eS)^{-1}.
\]
Both factors on the right are invertible integral matrices:
they reduce to \(U,U'\), respectively, modulo the maximal
ideal.  Hence \(\gamma\) preserves \(L_0\) at this prime and
has an integral inverse there.
We did not assume local integrality of \(\gamma\).

We may now reduce the same difference equation modulo
\(\mathfrak q^e\).
It gives \(U=\gamma U'\) in that quotient.
Using the \(\beta\)-basis and the definition of the two
labels, this says
\[
                  \gamma\beta k'=\beta k .
\]
This holds at every prime dividing \(\mathfrak a\).
The Chinese remainder theorem therefore gives one matrix
\[
                   k_\gamma=k(k')^{-1}\in J_{\rm lev}
\]
for the action of \(\gamma\) on
\(L_0/\mathfrak aL_0\).
Clearing denominators prime to \(\mathfrak a\) expresses the
same assertion as the source's \srcid{Level.Red}.
Thus \(\gamma\) is one of the generators defining \(\Delta\),
a contradiction.
\end{proof}

This proves the support condition \srcid{IsSystem.supp}.
Together with
Propositions~\ref{prop:sr-exact-level-gram}
and~\ref{prop:sr-level-isotropic-relation}, the common vacuum
multiplier, lattice-counting estimate and Neumann inverse,
the finite test-vector system now has explicit norm,
off-diagonal, automorphy, outside-level and isotropic
translation proofs in the same model.
The global doubled theta identity and the arithmetic
realization of that system still have to be matched to it.

\subsection{A nonzero-degree relation along every isotropic line}
\label{subsec:sr-level-isotropic-relation}

Use the fixed ambient model and the level test functions of
Section~\ref{sec:sr-level-phase-proof}.  Assume the level ideal is
proper.  Its conjugate is then proper as well.

\begin{srlemma}\label{lem:sr-level-line-detection}
Let \(0\ne\xi\in E^r\), and let \(y\) be any row tuple whose
first \(r\) rows span \(L_0/\mathfrak aL_0\).
For every allowed \(L,\Lambda\), there is a vector \(v\in E^n\)
such that
\[
 \delta=v\otimes\xi\in A_{\Lambda,L},\qquad
                         B_n(\delta,y)\notin\mathbb Z .
\]
Here the row tuple \(y\) can be any representative of a
nonzero level label.
\end{srlemma}

\begin{proof}
Put \(N=(E\xi)\cap\mathfrak jL\), an
\(\mathcal O_E\)-lattice of rank one on the line \(E\xi\).
It suffices to find \(\delta\in\Lambda\otimes N\) with the
stated nonintegral pairing, since every tensor in this module
has all its rows on \(E\xi\).
Localizing the intersection is legitimate: after bounding
denominators it is the kernel of a map of finitely generated
modules, and localization is exact.

Choose a prime \(\mathfrak q\mid\mathfrak a\), write
\(e=v_{\mathfrak q}(\mathfrak a)>0\), and put
\(\mathfrak p=\overline{\mathfrak q}\).
Coprimality of \(\mathfrak a,\overline{\mathfrak a}\)
gives \(\mathfrak p\ne\mathfrak q\), with
\[
 \mathfrak j_{\mathfrak p}
       =\mathfrak p^{-e}\mathcal O_{E,\mathfrak p},
 \qquad L_{\mathfrak p}=(L_0)_{\mathfrak p},\qquad
 \Lambda_{\mathfrak p}=\mathcal O_{E,\mathfrak p}^n .
\]
The last assertion holds at every prime over the underlying
rational prime \(\ell\): this prime divides \(M\), and
\(\Lambda\) is standard at \(M\).
Use completions in the rest of this local argument.

The saturated rank-one module
\(E_{\mathfrak p}\xi\cap(L_0)_{\mathfrak p}\)
has a primitive generator \(u\).
Indeed a generator lying in
\(\mathfrak p(L_0)_{\mathfrak p}\) could be divided by a
uniformizer and would not generate the intersection.
Consequently
\[
       N_{\mathfrak p}
            =\mathfrak p^{-e}\mathcal O_{E,\mathfrak p}u
                      \not\subset (L_0)_{\mathfrak p}.
\]
The trace pairing at \(\ell\) pairs the
\(\mathfrak p\)-component only with the
\(\mathfrak q\)-component.  This follows from
\(B(az,z')=B(z,\overline a z')\), applied to the idempotents
of \(E\otimes\mathbb Q_\ell\).
Selfduality of \(L_0\) says exactly that these two local
components are integral duals.  This formulation includes the
local different already present in the trace form.

Suppose every \(a\in N_{\mathfrak p}\) paired integrally with
each row \(y_i\), \(1\le i\le r\), at \(\mathfrak q\).
The same holds after multiplying each row by any local
integer, since the adjoint scalar preserves \(N_{\mathfrak p}\).
The rows span
\((L_0)_{\mathfrak q}/\mathfrak q^e(L_0)_{\mathfrak q}\).
They therefore express every \(z\in(L_0)_{\mathfrak q}\) as
a sum of such integer multiples plus an element of
\(\mathfrak q^e(L_0)_{\mathfrak q}\).
The remainder also pairs integrally with \(N_{\mathfrak p}\),
because
\(\mathfrak p^eN_{\mathfrak p}
 \subset(L_0)_{\mathfrak p}\).
It would follow that all of \(N_{\mathfrak p}\) pairs
integrally with \((L_0)_{\mathfrak q}\).
The local duality just established would give
\(N_{\mathfrak p}\subset(L_0)_{\mathfrak p}\), a contradiction.
Thus some row and some \(a\in N_{\mathfrak p}\) have
nonintegral \(\mathbb Q_\ell\)-pairing.

Place \(a\) in that row and zero in the other rows and prime
components of the \(\ell\)-adic completion of
\(\Lambda\otimes N\).
Global lattice points are dense in this completion: choose a
\(\mathbb Z\)-basis of the lattice and use the density of
\(\mathbb Z\) in \(\mathbb Z_\ell\) in each coordinate.
Approximate this local tensor closely enough that its pairing
with the fixed rational tuple \(y\) changes by an element of
\(\mathbb Z_\ell\).
The resulting global tensor \(\delta\) still has nonintegral
\(\ell\)-adic pairing, so \(B_n(\delta,y)\notin\mathbb Z\).
\end{proof}

\begin{srproposition}\label{prop:sr-level-isotropic-relation}
If in addition \(\xi\) is isotropic for the polarized
Hermitian form on \(E^r\), then there are a finite set
\(S\subset E^n\) and complex numbers \(a_v\) with
\[
 \sum_{v\in S}a_v\ne0,\qquad
 \sum_{v\in S}a_v S(v\otimes\xi)\Phi_{\Lambda,L}=0 .
\]
All these are translations in the same Schr\"odinger model.
\end{srproposition}

\begin{proof}
Let
\[
 {\cal B}_\xi=\{v\in E^n:v\otimes\xi\in A_{\Lambda,L}\}.
\]
The Hermitian isotropy gives
\(B(a\xi,b\xi)=0\) for every \(a,b\in E\).
It follows that \(S(v\otimes\xi)\), \(v\in{\cal B}_\xi\),
form an ordinary additive representation, with no Heisenberg
cocycle.  Also
\(\epsilon_w(v\otimes\xi)=1\), so these operators fix the
vacuum \(v_{\Lambda,L}\).

For a level representative \(y_k\), commuting a translation
past \(S(y_k)\) gives
\[
 S(v\otimes\xi)S(y_k)v_{\Lambda,L}
   =\chi_k(v)S(y_k)v_{\Lambda,L},\qquad
 \chi_k(v)=\exp\{-2\pi iB_n(v\otimes\xi,y_k)\}.
\]
This is a character of the additive group
\({\cal B}_\xi\).
Lemma~\ref{lem:sr-level-line-detection} supplies
\(v_k\in{\cal B}_\xi\) with \(\chi_k(v_k)\ne1\)
for each of the finitely many labels.
In the complex group algebra of \({\cal B}_\xi\) take
\[
 Z_\xi=\prod_{k\in J_{\rm lev}}
                   ([v_k]-\chi_k(v_k)[0]).
\]
Evaluation under \(\chi_k\) kills this product, by its
\(k\)-th factor.  Thus its translation action kills every
\(S(y_k)v_{\Lambda,L}\), and hence their weighted sum
\(\Phi_{\Lambda,L}\).
Its augmentation is
\[
       \prod_{k\in J_{\rm lev}}(1-\chi_k(v_k))\ne0 .
\]
Writing out the finite product as
\(Z_\xi=\sum_{v\in S}a_v[v]\) proves the assertion.
\end{proof}

This proves the finite test-function relation called
\srcid{IsSystem.cusp}, including the required nonzero
coefficient sum.  Passing this relation through an absolutely
convergent invariant integral makes that integral zero.
Using it to eliminate every nongraph term in a doubled
Eisenstein expansion still requires that expansion, its
identification of the corresponding isotropic stabilizer,
and the convergence or regularization justified for that
term.  Those analytic identifications are not consequences
of the finite group-algebra calculation alone.

\subsection{Invariant heat jets from the scalar Siegel--Weil identity}
\label{subsec:sr-invariant-heat-jets}

The scalar identity at zero elliptic variable does not, by
itself, give a Jacobi identity at arbitrary elliptic variable.
The following algebraic and analytic argument explains the
invariant differentiations actually needed for the theta
kernel.

Let \(T\) be a finite set, \(n\ge r\), and for each
\(\tau\in T\) write
\[
 U_\tau\in\operatorname{Mat}_{n,r}(\mathbb C),\qquad
 V_\tau\in\operatorname{Mat}_{r,n}(\mathbb C).
\]
Use the action
\(U_\tau\mapsto g_\tau U_\tau\),
\(V_\tau\mapsto V_\tau g_\tau^{-1}\), and put
\(s_{\tau,ab}=\sum_k(U_\tau)_{ka}(V_\tau)_{bk}\).

\begin{srlemma}\label{lem:sr-invariants-pairings}
Every complex polynomial invariant under all
\(g_\tau\in U(n)\) is a polynomial in the
\(s_{\tau,ab}\).
\end{srlemma}

\begin{proof}
Invariance under the unitary group implies invariance under
\(\operatorname{GL}_n(\mathbb C)\).
Indeed differentiating the invariant polynomial identity
along \(\exp(tK)\), for every skew-Hermitian \(K\), gives
zero infinitesimal action.
The complex span of these \(K\)'s is the full matrix algebra:
every matrix is the sum of a skew-Hermitian matrix and
\(i\) times a skew-Hermitian matrix.
The same infinitesimal identity therefore holds for every
complex matrix.  It gives invariance along its exponential.
Every invertible complex matrix has a logarithm, by putting
it in Jordan form and using the finite logarithm series on
each nilpotent part, so the full general linear group acts
invariantly.

For one \(\tau\), put
\(U_0=\binom{I_r}{0}\).
On the open set where \(U\) has rank \(r\), a general linear
change of rows carries \(U\) to \(U_0\).
Write the corresponding \(V\) as \([C,D]\), so \(C=VU\).
On the further open set where \(C\) is invertible, the
stabilizer of \(U_0\) contains
\[
       g=\begin{pmatrix}I_r&A\\0&I_{n-r}\end{pmatrix},
                   \qquad A=C^{-1}D.
\]
Right multiplication by \(g^{-1}\) sends \([C,D]\) to
\([C,0]\).  Thus an invariant polynomial \(F\) equals
\(F(U_0,[VU,0])\) on this open set.
The last expression is a polynomial in \(VU\).
The open set is nonempty and Zariski dense, so the polynomial
identity holds everywhere.
For \(n=r\) the same proof omits the empty lower block.
Apply this independently at each \(\tau\); the resulting
polynomial is obtained by fixing all \(U_\tau=U_0\) and
putting \(V_\tau=[C_\tau,0]\).
\end{proof}

For elliptic coordinates \(u_{\tau,ka},v_{\tau,bk}\), define
the constant-coefficient contraction operators
\[
 {\cal D}_{\tau,ab}
       =\sum_{k=1}^n
          \frac{\partial^2}
           {\partial u_{\tau,ka}\,\partial v_{\tau,bk}} .
\]
Let \(Z\) vary in an open complex matrix domain \(\Omega\).
Suppose \(D(Z,u,v)\) is holomorphic on
\(\Omega\times\mathbb C^{2|T|nr}\) and satisfies
\begin{equation}\label{eq:sr-exact-heat}
 {\cal D}_{\tau,ab}D
         =2\pi i\,\frac{\partial D}{\partial Z_{\tau,ab}},
                 \qquad D(Z,0,0)=0 .
\end{equation}
The coordinate order in \(Z_{\tau,ab}\) is fixed by this
display.  For a theta exponential
\[
 \exp\!\left(2\pi i\sum_k
       [\,x_kZ\overline{x_k}^{\,t}+x_ku_k+
                                      v_k\overline{x_k}^{\,t}\,]\right),
\]
direct differentiation gives precisely this heat equation,
with the corresponding row and column convention.

\begin{srproposition}\label{prop:sr-heat-invariant-zero}
For every invariant polynomial \(Q\) in the differential
symbols,
\[
                 Q(\partial_u,\partial_v)
                                    D(Z,0,0)=0 .
\]
\end{srproposition}

\begin{proof}
Lemma~\ref{lem:sr-invariants-pairings} writes
\(Q=p(s_{\tau,ab})\).
Substituting the commuting derivatives gives
\(Q(\partial_u,\partial_v)=p({\cal D}_{\tau,ab})\).
All the derivatives commute on holomorphic functions.
Repeated use of \eqref{eq:sr-exact-heat} therefore yields
\[
 p({\cal D}_{\tau,ab})D(Z,0,0)
   =p(2\pi i\,\partial_{Z_{\tau,ab}})D(Z,0,0)=0.
\]
The last equality follows by differentiating the identically
zero holomorphic function \(Z\mapsto D(Z,0,0)\).
\end{proof}

A Gaussian multiplier can be absorbed without changing this
conclusion.  Let
\[
 G(u,v)=c\exp\!\left(\sum_{\tau,k,a,b}
                     A_{\tau,ab}u_{\tau,ka}v_{\tau,bk}\right).
\]
It is invariant under the row action above.
For a polynomial differential operator \(Q\), Leibniz's rule
gives a unique polynomial \(p_G\) such that
\begin{equation}\label{eq:sr-absorb-gaussian}
   Q(\partial)(G F)(0)=p_G(\partial)F(0)
\end{equation}
for all holomorphic \(F\).
For example it is specified explicitly by
\[
              p_G(\xi)
                  =\left.Q(\partial_x)
                       \bigl(G(x)e^{\langle x,\xi\rangle}\bigr)
                                                     \right|_{x=0}.
\]
The right side is polynomial in \(\xi\).
Coefficient comparison on monomials, or direct Leibniz
expansion, proves \eqref{eq:sr-absorb-gaussian} for every
holomorphic \(F\); only finitely many derivatives occur.
If \(Q\) and \(G\) are invariant, the chain rule in this
exponential formula shows that \(p_G\) is invariant under
the dual row action.
Applying Proposition~\ref{prop:sr-heat-invariant-zero} to
\(p_G\) proves
\[
                   Q(\partial)(G D)(Z,0,0)=0 .
\]

The same argument permits fixed linear column substitutions.
More precisely, if the two sets of differentiating
directions intertwine the row action, their pullback of an
invariant bidifferential is invariant.
Right multiplication by the fixed matrices in the doubled
coordinates commutes with the left row action; the dual
variables carry the inverse action.
For a family \(P_j\) satisfying the common invariant-sum
identity, the pulled-back operator
\(\sum_jP_j(\partial_z)\overline{P_j}(\partial_w)\)
therefore has the required invariance.
The Gaussian factor is of the displayed row-contraction
form, so it can be absorbed as above.
This proves the invariant-jet implication used in the
doubled frame-kernel construction, once its stated column
coordinates are used in the same order.

In particular, to compare the class-averaged differentiated
theta series with differentiated Eisenstein series, it is
enough to prove the scalar identity at \((u,v)=(0,0)\),
the exact heat equation, local holomorphy, and normal
convergence permitting the stated differentiations.
One need not assert a stronger Jacobi identity.
This proposition proves that implication; it does not
supply the scalar Siegel--Weil identity or its
class-averaging normalization.
